\documentclass[10pt,aspectratio=169]{beamer}

\usetheme{metropolis}
\metroset{progressbar=frametitle}

\usepackage[T1]{fontenc}
\usepackage[utf8]{inputenc}
\usepackage[brazilian]{babel}
\usepackage{booktabs}
\usepackage{listings}
\usepackage{xcolor}

\definecolor{bgcode}{RGB}{248,248,242}
\definecolor{kwcode}{RGB}{36,82,160}
\definecolor{strcode}{RGB}{120,40,140}
\definecolor{cmtcode}{RGB}{60,120,60}
\definecolor{numcode}{RGB}{130,130,130}
\definecolor{bgterm}{RGB}{30,30,30}
\definecolor{fgterm}{RGB}{230,230,230}

\lstdefinestyle{cstyle}{
    language=C,
    backgroundcolor=\color{bgcode},
    commentstyle=\color{cmtcode}\itshape,
    keywordstyle=\color{kwcode}\bfseries,
    numberstyle=\tiny\color{numcode},
    stringstyle=\color{strcode},
    basicstyle=\ttfamily\scriptsize,
    breaklines=true,
    keepspaces=true,
    numbers=left,
    numbersep=5pt,
    showstringspaces=false,
    tabsize=4,
    frame=single,
    framerule=0pt,
    morekeywords={uint8_t,uint32_t,size_t,__ESBMC_assert}
}

\lstdefinestyle{outstyle}{
    backgroundcolor=\color{bgterm},
    basicstyle=\ttfamily\tiny\color{fgterm},
    breaklines=true,
    keepspaces=true,
    showstringspaces=false,
    tabsize=4,
    frame=single,
    framerule=0pt
}

\title{Operational Models para Verifica\c{c}\~ao de Brokers MQTT com ESBMC}
\subtitle{Defesa Intermedi\'aria}
\author{Gustavo Lima \and Willian Jean \and Paulo Gomes}
\date{Junho de 2026}

\begin{document}

\maketitle

% ------------------------------------------------------------------
\begin{frame}{Roteiro}
    \tableofcontents[hideallsubsections]
\end{frame}

% ==================================================================
\section{Progresso}

\begin{frame}{O que foi feito / O que \'e novo}
    \begin{columns}[T]
        \column{0.48\textwidth}
        \textbf{Estabelecido anteriormente}
        \begin{itemize}
            \item ESBMC com BMC para C
            \item Conceito de harness simb\'olico
            \item 7 CVEs do Mosquitto modelados
            \item Problema: APIs de rede sem modelo
            \item Modelos: \texttt{socket}, \texttt{recv}, \texttt{send}, \texttt{bind}, \texttt{accept}, \texttt{close}
        \end{itemize}
        \column{0.48\textwidth}
        \alert{\textbf{Desde ent\~ao}}
        \begin{itemize}
            \item Modelos \texttt{select}/\texttt{poll} conclu\'idos
            \item 4 harnesses de integra\c{c}\~ao via \texttt{recv()}
            \item \textbf{21 harnesses totais verificados}
            \item \alert{Nova vulnerabilidade encontrada no Mosquitto~2.1.x (CWE-191)}
        \end{itemize}
    \end{columns}
\end{frame}

\begin{frame}{Modelos Operacionais Implementados}
    \textbf{Arquivo:} \texttt{esbmc\_models/network\_stubs.c}

    \vspace{0.5em}
    \begin{tabular}{ll}
        \toprule
        Fun\c{c}\~ao & Modelo \\
        \midrule
        \texttt{socket} & fd nondet em $[3, 63]$ \\
        \texttt{bind}, \texttt{listen} & retornam 0 \\
        \texttt{accept} & fd cliente nondet \\
        \texttt{recv} & preenche buffer com bytes nondet \\
        \texttt{send} & retorna comprimento nondet \\
        \texttt{close} & retorna 0 \\
        \texttt{select} \textbf{(novo)} & seta bits nondet em \texttt{fd\_set} \\
        \texttt{poll} \textbf{(novo)} & preenche \texttt{revents} nondet \\
        \bottomrule
    \end{tabular}

    \vspace{0.4em}
    \small Os novos modelos \texttt{select}/\texttt{poll} cobrem o loop de
    eventos do broker, antes inacess\'ivel \`a verifica\c{c}\~ao.
\end{frame}

% ==================================================================
\section{Resultados}

\begin{frame}{CVEs Cobertos (or\'aculo de valida\c{c}\~ao)}
    \small
    \begin{tabular}{p{1.9cm}p{3.2cm}lc}
        \toprule
        CVE & Tipo & CWE & Fix \\
        \midrule
        2017-7651 & Integer Overflow & 190 & 1.4.15 \\
        2017-7654 & Memory Leak      & 401 & 1.4.15 \\
        2017-7650 & ACL Bypass       & 284 & 1.4.12 \\
        2023-3592  & Memory Leak     & 401 & 2.0.16 \\
        2023-0809  & Excess Alloc    & 400 & 2.0.16 \\
        2023-28366 & Memory Leak QoS2& 401 & 2.0.16 \\
        2024-3935  & Double Free     & 415 & 2.0.19 \\
        \bottomrule
    \end{tabular}

    \vspace{0.3em}
    + 2 harnesses de propriedade MQTT~v5 \quad
    + 4 harnesses de integra\c{c}\~ao via \texttt{recv()}
\end{frame}

\begin{frame}{Resultados: Detec\c{c}\~ao}
    \scriptsize
    \begin{tabular}{p{3.0cm}lcc}
        \toprule
        Subsistema & Flags & Vuln. & Fix \\
        \midrule
        \multicolumn{4}{l}{\textit{CVE -- isolados}} \\
        \texttt{remaining\_len}    & uof        & \alert{FAIL} & -- \\
        \texttt{connect\_memleak}  & mlc        & \alert{FAIL} & -- \\
        \texttt{acl\_bypass}       & uw~32      & \alert{FAIL} & SUCC \\
        \texttt{will\_properties}  & mlc        & \alert{FAIL} & SUCC \\
        \texttt{initial\_packet}   & uw~4       & \alert{FAIL} & SUCC \\
        \texttt{qos2\_dup}         & mlc        & \alert{FAIL} & SUCC \\
        \texttt{bridge\_remap}     & mlc+uw~16  & \alert{FAIL} & SUCC \\
        \multicolumn{4}{l}{\textit{Integra\c{c}\~ao / Propriedades}} \\
        \texttt{pkt\_recv\_*} (2)  & net\_stubs & \alert{FAIL} & SUCC \\
        \texttt{varint}/\texttt{max\_pkt}     & uw~8       & --           & SUCC \\
        \multicolumn{4}{l}{\textit{Nova falha (2.1.x)}} \\
        \texttt{proxy\_v2\_tlv}   & uof+assert & \alert{FAIL} & SUCC \\
        \bottomrule
    \end{tabular}
    \vspace{0.2em}\\
    {\tiny uof=\texttt{--unsigned-overflow-check}; mlc=\texttt{--memory-leak-check}; uw=\texttt{--unwind}; assert=\texttt{\_\_ESBMC\_assert} manual}
\end{frame}

% ==================================================================
\section{Nova Vulnerabilidade (Mosquitto 2.1.x)}

\begin{frame}{Resultados: Desempenho (ESBMC~8.3, z3)}
    \small
    \begin{tabular}{lcrr}
        \toprule
        Harness & Resultado & Tempo (s) & Mem.\ (MB) \\
        \midrule
        \texttt{remaining\_length} (vuln.)   & \alert{FAILED}    &   1,0 &   97 \\
        \texttt{connect\_memleak} (vuln.)    & \alert{FAILED}    &   1,5 &  166 \\
        \texttt{acl\_bypass} (unwind 32)     & \alert{FAILED}    &  87,6 & 2220 \\
        \texttt{acl\_bypass\_fixed}          & SUCCESSFUL        &  23,1 & 1996 \\
        \texttt{will\_properties} (vuln.)    & \alert{FAILED}    &   1,0 &   98 \\
        \texttt{qos2\_dup} (vuln.)           & \alert{FAILED}    &   1,3 &  119 \\
        \texttt{bridge\_remap} (vuln.)       & \alert{FAILED}    &   1,0 &   78 \\
        \texttt{pkt\_recv\_rem\_len} (vuln.) & \alert{FAILED}    &   0,6 &   77 \\
        \texttt{proxy\_v2\_tlv} (vuln., 2.1.x) & \alert{FAILED} &   0,9 &   73 \\
        \texttt{proxy\_v2\_tlv} (fixed)      & SUCCESSFUL        &   0,5 &   73 \\
        \emph{demais harnesses}              & SUCCESSFUL        & $<4$  & $<$135 \\
        \bottomrule
    \end{tabular}

    \vspace{0.4em}
    \small 19/21 harnesses concluem em $<$~4\,s e $<$~170\,MB.
    \texttt{acl\_bypass} \'e outlier (unwind~32).
\end{frame}

\begin{frame}{Vulnerabilidade Encontrada no Mosquitto~2.1.x (1/2)}
    \begin{block}{Falha em \texttt{src/proxy\_v2.c} -- \texttt{read\_tlv\_ssl} (CWE-191)}
        Parser de sub-TLVs SSL no PROXY Protocol~v2.\\[0.3em]
        \textbf{Bug:} \texttt{len = (uint16\_t)(len - (3 + tlv\_len))} executado\\sem verificar
        \texttt{3 + tlv\_len \textbf{<=} len} antes.
    \end{block}
    \vspace{0.3em}
    \textbf{Contraexemplo ESBMC ($<$1\,s):}
    \small
    \begin{tabular}{ll}
        \texttt{ssl\_tlv\_len} = 256 & \texttt{len} = 251 (ap\'os 5-byte SSL header) \\
        \texttt{tlv\_len} = 253       & \textit{passa} no buffer check ($\leq$492) \\
        \multicolumn{2}{l}{\texttt{3+253=256 > 251} \;\textbf{$\Rightarrow$ len wraps para 65.531}}
    \end{tabular}
    \vspace{0.3em}\\
    \alert{FAILED} (vuln.) \; / \; SUCCESSFUL (corre\c{c}\~ao: adicionar guarda)\\[0.2em]
    \footnotesize Requer \texttt{proxy\_protocol~true}. Ser\'a reportado ao Eclipse (\emph{coordinated disclosure}).
\end{frame}

\begin{frame}{Vulnerabilidade Encontrada no Mosquitto~2.1.x (2/2)}
    \begin{alertblock}{Por que nem flags autom\'aticas nem fuzzer detectaram}
        \small
        A subtra\c{c}\~ao \texttt{len - (3+tlv\_len)} ocorre em aritm\'etica \texttt{int}
        com sinal (promo\c{c}\~ao impl\'icita do C).\\[0.3em]
        \begin{tabular}{ll}
            \texttt{--unsigned-overflow-check} & n\~ao acionado (opera\c{c}\~ao \'e em \texttt{int}) \\
            \texttt{--overflow-check}          & $-4$ \'e \texttt{int} v\'alido \\
            UBSan / fuzzer (OSS-Fuzz)          & sem \emph{crash}, comportamento definido em C \\
        \end{tabular}
    \end{alertblock}
    \vspace{0.3em}
    \begin{block}{Por que o ESBMC com harness detectou}
        \small
        O harness especifica o \textbf{invariante de dom\'{\i}nio} como asser\c{c}\~ao:\\[0.2em]
        \texttt{\_\_ESBMC\_assert((uint32\_t)3 + tlv\_len <= (uint32\_t)len, ...)}\\[0.2em]
        \'E exatamente esse tipo de propriedade --- n\~ao um crash ---
        que a verifica\c{c}\~ao formal com asser\c{c}\~ao de dom\'{\i}nio consegue verificar.
    \end{block}
    \vspace{0.2em}
    \footnotesize Commit \texttt{23c918ee} (jan./2026): fuzzer para \texttt{proxy\_v2} adicionado. Este bug n\~ao foi encontrado pelo fuzzer.
\end{frame}

% ==================================================================
\section{Pr\'oximos Passos e Cronograma}

\begin{frame}{Pr\'oximos Passos}
    \begin{columns}[T]
        \column{0.48\textwidth}
        \textbf{Limita\c{c}\~oes atuais}
        \begin{itemize}\small
            \item An\'alise \emph{context-bounded}
            \item CVE-2024-8376 ainda n\~ao modelado
            \item APIs pendentes: \texttt{epoll\_*}, TLS/SSL
            \item Modelos externos ao \texttt{c2goto}
        \end{itemize}
        \column{0.48\textwidth}
        \textbf{A fazer}
        \begin{enumerate}\small
            \item Divulgar falha ao Eclipse (\emph{disclosure})
            \item Integrar modelos ao \texttt{c2goto}
            \item Buscar mais falhas no Mosquitto 2.1.x
            \item Harness para CVE-2024-8376
            \item Submeter artigo final
        \end{enumerate}
    \end{columns}
\end{frame}

\begin{frame}{Cronograma do trabalho}
    \small
    \begin{tabular}{clcc}
        \toprule
        \# & Atividade & Per\'iodo & Status \\
        \midrule
        1 & Pesquisa bibliogr\'afica e CVEs       & Mar--Abr/26 & \checkmark \\
        2 & Harnesses isolados                    & Mar--Mai/26 & \checkmark \\
        3 & Operational models (socket/recv/send) & Abr--Mai/26 & \checkmark \\
        4 & Modelos select/poll + integra\c{c}\~ao     & Mai/26      & \checkmark \\
        5 & Experimentos e m\'etricas              & Mai--Jun/26 & \checkmark \\
        6 & Defesa intermedi\'aria                 & Jun/26      & \textbf{agora} \\
        7 & Integra\c{c}\~ao c2goto + defesa final     & Jul/26      & pendente \\
        \bottomrule
    \end{tabular}
\end{frame}

% ==================================================================
\begin{frame}[standout]
    Obrigado.

    \vspace{0.8em}
    \small Reposit\'orio: \texttt{harnesses/mosquitto/} \quad
    Logs: \texttt{*.log} \quad
    Modelos: \texttt{esbmc\_models/}
\end{frame}

\end{document}